\section*{第 \ref{ch:g9:energy-conversions} 章 --- 能量的形式与转化}

\begin{solution}{exo:g9:energy-conversions:1}
动能（滚动的小球）、势能（提起的水桶）、
热能（温热的洗澡水）、化学能（早餐、电池）、
电能（墙上的供电）、辐射能（穿过
窗户的阳光）、核能（发电站的堆芯）。
\end{solution}

\begin{solution}{exo:g9:energy-conversions:2}
水桶：$2 \times 9.8 \times 10 = \qty{196}{J}$。登山者：$60
\times 9.8 \times 500 = \qty{2.9e5}{J}$ --- 接近
三百千焦的重力势能。
\end{solution}

\begin{solution}{exo:g9:energy-conversions:3}
在每一过程中，能量改变形式与归属，而总量
恰好守恒。永动机正是死于此：
输出多于其所储存或接收能量的机器，
会凭空产生能量 --- 绝无申辩余地。
\end{solution}

\begin{solution}{exo:g9:energy-conversions:4}
摆至最高点：全为势能，瞬时静止。最低点：全为
动能，速率最大。能量流失：每次经过时，都有少许
转化为空气与支点的热能，秋千逐渐停下 --- 能量
只是转移，总量未变。
\end{solution}

\begin{solution}{exo:g9:energy-conversions:5}
$v = \sqrt{2 \times 9.8 \times 5} \approx \qty{9.9}{m/s}$ ---
约 \qty{36}{km/h}。令 $m g h$ 与
$\frac12 m v^2$ 相等时，质量相消：轻重跳水者
入水速率相同。
\end{solution}

\begin{solution}{exo:g9:energy-conversions:6}
有用能量与总能量之比：$170 \div 200 = 0.85$ --- 百分之
八十五。缺失的 \qty{30}{J} 转为绕组与
轴承中的热能。
\end{solution}

\begin{solution}{exo:g9:energy-conversions:7}
输入能量中只有百分之五以光的形式输出，其余百分之
九十五以热能形式散失。按输出归类，它其实是个暖炉 ---
只不过恰好会发光。
\end{solution}

\begin{solution}{exo:g9:energy-conversions:8}
化学能（早餐）$\to$ 肌肉做功 $\to$ 骑手与车的
动能 $\to$ 以势能 $m g
h$ 的形式储存于山顶，沿途处处有热能泄漏：温热的肌肉、
被搅动的暖空气、发热的轮胎。总量守恒。
\end{solution}

\begin{solution}{exo:g9:energy-conversions:9}
每吨：$1000 \times 9.8 \times 50 = \qty{4.9e5}{J}$；
按百分之九十计，约为 \qty{4.4e5}{J} 的电能。对于
\qty{4.4e6}{W} 的功率：$\num{4.4e6} \div \num{4.4e5} = 10$ 吨
每秒流经涡轮机。
\end{solution}

\begin{solution}{exo:g9:energy-conversions:10}
碰撞前：$0.5 \times 9.8 \times 2 = \qty{9.8}{J}$。
弹起后：$0.5 \times 9.8 \times 1.2 \approx \qty{5.9}{J}$。
损耗：约 \qty{3.9}{J} 转化为小球与地面的热能（以及一
声轻响，其声音本身也很快变成热能）。核算
在回弹高度降为零时结束：原先的 $9.8$ 焦耳
最终全部计入热能账户。
\end{solution}

\begin{solution}{exo:g9:energy-conversions:11}
$E_p = 500 \times 9.8 \times 30 = \qty{1.47e5}{J}$；效率为百分之九十时，
电表须提供 $\num{1.47e5} \div 0.9 \approx
\qty{1.6e5}{J}$。多出的十分之一并未
流失到宇宙之外 --- 它使电动机、齿轮与转轴发热；
能量守恒仍然成立，升降机的账目中只是多了一项热能。
\end{solution}

\begin{solution}{exo:g9:energy-conversions:12}
最低点速率：$v = \sqrt{2 \times 9.8 \times 4} \approx
\qty{8.9}{m/s}$。理想情况下，滑板者应恰好回到
\qty{4}{m} 高的右侧边缘；实际上会略低一些，每次
经过都更低一点，直到滑板者在底部摆动至静止
--- 整个 $m g h$ 最终化作车轮、坡道的热能，
以及空气中的热能。能量并未消灭，但有用的部分已被耗尽：
世界的账目多了分散的热能，而
滑板者失去了集中、可再用的高度 ---
能量的数量犹存，但其用处再不能复原。
\end{solution}

\begin{solution}{pb:g9:energy-conversions:1}
\textbf{1.} $E_p = 2000 \times 9.8 \times 45 =
\qty{8.8e5}{J}$。
\textbf{2.} $v = \sqrt{2 \times 9.8 \times 45} \approx
\qty{29.7}{m/s} \approx \qty{107}{km/h}$。
\textbf{3.} 高 \qty{50}{m} 的峰顶所需能量超过了
最高点所给出的能量：账目不允许。从 \qty{45}{m} 高处
静止出发，此后任何山丘都不得超过 \qty{45}{m} ---
实际上，摩擦使其明显更低。
\textbf{4.} 实测值：$0.5 \times 2000 \times 28^2 =
\qty{7.8e5}{J}$，对比理想值 \qty{8.8e5}{J}：约
\qty{1e5}{J} 计入热能账户 --- 轨道、
车轮与空气发热了。
\textbf{5.} $P = \num{8.8e5} \div 60 \approx
\qty{1.5e4}{W}$ --- 约十五千瓦。
\textbf{6.} $14.7 \div 20 \approx 0.74$：百分之
七十四。
\textbf{7.} 十小时内，每四分钟一次，共 $150$
次升降；每次耗电 $20 \times \tfrac{1}{60} =
\qty{0.33}{kWh}$；每日 $150 \times 0.33 = \qty{50}{kWh}$
--- $12.5$ 个货币单位。
\textbf{8.} 一台性能优良的制动发电机或许能回收相当可观的
到达时的动能 --- 但每个环节
（轨道、制动器、发电机、变压器）都要缴纳热能
的“手续费”，回收的焦耳永远无法等于绞车的
耗能：能量守恒允许循环利用；
效率小于一，不允许完整闭环。这趟游乐
只能部分抵偿自身消耗，却永远无法自给自足。
\textbf{9.} $v = \sqrt{2 \times 9.8 \times 50} \approx
\qty{31.3}{m/s}$。
\textbf{10.} 入口处动能 $0.5 \times 500 \times 980 =
\qty{2.45e5}{J}$，加上最后十米的 $500 \times 9.8
\times 10 = \qty{4.9e4}{J}$：约 \qty{2.9e5}{J} 须由
制动器承担 --- 经感应涡流，转为
制动片的热能：感应，就是一份“退休计划”。
\textbf{11.} 该说法成立，措辞虽未必精确，
但感觉上确是如此：吊篮与乘客一同下落，
同步下落的人彼此不受压力 --- 尖叫声源自共同的自由落体，
重力完整地存在，却无从察觉。
\textbf{12.} 例如：“我们的过山车只做一笔买卖：
高度换速度，再换回高度，即 $m g h$ 与 $\frac12 m v^2$ 之间的交换。
每一座游乐设施都遵循同一条定律：能量总量，
在每一点核算，都从未改变。还有一本账向它们全部征收：
热能 --- 摩擦悄悄收取的费用，
每次乘坐都要缴纳。”
\end{solution}
